\chapter{Fusion and Radioactive Decay}\label{ch:nuclear-decay}
\index{fusion!solar}\index{fusion!muon-catalyzed}\index{alpha decay}\index{beta decay}\index{Geiger--Nuttall law}\index{neutrino}

\begin{quote}
\itshape A model earns trust by predicting where it works and saying plainly where it breaks.\\
\normalfont\hfill --- working note, 2026.
\end{quote}
\bigskip

If the nucleus is bound by Coulomb confinement (Chapter~\ref{ch:nuclear-binding}), then the processes that make
and unmake nuclei --- fusion and radioactive decay --- should follow from the same energetics, with no separate
strong or weak force. This chapter reports four such computations. Two succeed cleanly (solar fusion energetics,
alpha-decay rates); one is a scaling result with clear caveats (muon-catalyzed fusion); and one \emph{partly
fails}, in an instructive way that we state without softening (beta decay).

\section{Solar fusion by Coulomb alone}\label{sec:solar-fusion}

The energy of the Sun is, on the standard account, the fusion of hydrogen to helium through the $pp$ chain,
$4p\to{}^4\mathrm{He}+2e^+ +2\nu_e$, converting two protons to neutrons and emitting positrons and neutrinos. In
the RealQM reading the ingredients are a hot, ionised soup of protons \emph{and electrons}, and the same net
result is reached by \emph{Coulomb rearrangement} without converting any proton into a neutron:
\[
p+e+p\;\to\;d, \qquad d+d\;\to\;{}^4\mathrm{He},
\]
and overall $4p+4e\to{}^4\mathrm{He}+2e$ --- four protons and four electrons in, a helium nucleus (two protons,
two confined electrons) plus two electrons out as the atomic shell. The released energy is the same
$\sim 26.7$~MeV, a state-function difference independent of path. The crux is the \emph{lepton bookkeeping}: no
lepton is created or destroyed in this path, so the process \emph{as written} does not require the emission of
positrons or neutrinos. We are careful here: this says the neutrino is not \emph{demanded} by the Coulomb path,
not that neutrinos do not exist; it is a possible Coulomb process, offered beside the standard one, not a
refutation of neutrino physics.

\section{Muon-catalyzed fusion as a change of glue}\label{sec:muon-fusion}

Muon-catalyzed fusion is, in this picture, a change of the binding \emph{glue}. Replace the electron in the
molecular ion $\mathrm{D}_2^+$ by a muon, $207$ times heavier, and the whole bound structure scales down by that
factor: the equilibrium separation $R_e\approx 2.0\,a_0$ of the electronic ion becomes $R_e/207\approx 512$~fm
for the muonic one, compressing the two deuterons by a factor of order $144$ relative to their free separation
and bringing them to distances where fusion proceeds. The RealQM computation reproduces this by rescaling the
$\mathrm{D}_2^+$ binding curve with the muon mass, and the result is a clean statement about \emph{scaling}: the
same Coulomb bound state, with a heavier glue, is a smaller bound state. The caveats are honest --- the curve is
computed in the Born--Oppenheimer approximation on a finite grid, and the per-domain mass that distinguishes
electron from muon is a route the solver supports but which is not a full dynamical treatment of the three-body
$dd\mu$ system. It is a scaling insight, not a rate calculation.

\section{Alpha decay: Coulomb tunnelling that works}\label{sec:alpha-decay}

Alpha decay is the escape of a He-4 unit through the Coulomb barrier of the residual nucleus --- Gamow's 1928
picture, and one that fits the Coulomb-confinement model without addition. The tunnelling probability through the
barrier, combined with the assault frequency, gives a decay rate that spans an enormous range as the decay energy
varies. Computed across the measured alpha emitters, the model reproduces the \emph{Geiger--Nuttall law} --- the
linear relation between the logarithm of the half-life and the inverse square root of the decay energy --- across
about \emph{twenty-four orders of magnitude} in half-life. This is the cleanest quantitative success of the
nuclear program: a parameter-light Coulomb-barrier calculation tracking rates from microseconds to billions of
years. The same tunnelling reach extends, more roughly, to cluster decay and to fission as the escaping unit
grows.

\section{Beta decay: a partial success, honestly}\label{sec:beta-decay}

Beta decay is where the program is tested hardest, and where it \emph{partly fails} --- reported here plainly,
because a model is only as trustworthy as its account of its own limits. Two things were attempted. First, a
deterministic \emph{phase-coincidence} model of the decay \emph{statistics}: treating the trigger as the
coincidence of internal phases, one recovers exponential decay and a sensible half-life distribution. That part
\emph{survives}. Second, and more ambitious, an attempt to dispense with the neutrino by having the available
energy carried entirely by the electron and a Coulomb field rearrangement. That part \emph{does not survive}: the
computed field is roughly two orders of magnitude too weak to carry the energy and momentum the neutrino accounts
for. The honest conclusion is therefore mixed and must be stated as such: \textbf{the deterministic
phase-coincidence account of decay timing stands; the elimination of the neutrino does not.} The neutrino is
neither removed nor explained away by this model; the three-body kinematics that motivated Pauli's neutrino in
1930 reasserts itself. We record the failure with the success, because the discipline of the whole program is to
let the computation, not the ambition, have the last word.
